\chapter{Manual técnico}

\section{Estructura relevante}

\begin{verbatim}
phase2-juan/
  simlab/
    environment.py      dominio de simulacion
    spec.py             validacion de entornos
    orchestrator.py     chat, herramientas y coordinacion
    architect.py        generacion de entorno
    tracker.py          observacion de simulaciones
    analyst.py          analisis de patrones
    reporter.py         generacion de informes PDF
    model_loader.py     carga de modelos de fase 1
    knowledge/          escritura de memorias de simulacion
    recall/             lectura del Knowledge Backbone
    api.py              FastAPI + WebSocket
  frontend/             cliente React
  tests/                pruebas pytest

shared/                 servicios comunes y migraciones
docker-compose.yml      PostgreSQL, MinIO, Neo4j y Qdrant
\end{verbatim}

\section{Servicios}

La ejecución completa requiere PostgreSQL, MinIO, Neo4j y Qdrant. El
acceso a modelos de lenguaje se realiza mediante OpenRouter. Voyage AI y
ZeroEntropy son opcionales: si no están configurados, el sistema puede
ejecutarse, pero la recuperación del Knowledge Backbone pierde calidad.

\section{Configuración}

La configuración vive en \texttt{phase2-juan/.env}. Las variables
principales son las conexiones a PostgreSQL, MinIO, Neo4j y Qdrant,
\texttt{OPENROUTER\_API\_KEY}, y los conmutadores
\texttt{ENABLE\_KNOWLEDGE\_READ} y \texttt{ENABLE\_KNOWLEDGE\_WRITE}.
El archivo \texttt{.env.example} contiene la lista completa.

\section{Comandos}

\begin{verbatim}
# Infraestructura para levantar los contenedores
docker compose up -d

# Backend
cd phase2-juan
uv sync
uv run uvicorn simlab.api:app --port 8000

# Frontend
cd phase2-juan/frontend
npm install
npm run dev

# CLI
cd phase2-juan && uv run simlab
\end{verbatim}

\section{Pruebas}

\begin{verbatim}
cd phase2-juan && uv run pytest tests/ -v
cd phase2-juan/frontend && npx playwright test
\end{verbatim}

\section{Extensión}

Un nuevo modelo de decisión generado por la primera fase queda
disponible para esta fase sin necesidad de tocar el código del
laboratorio: la primera fase lo registra en la tabla \texttt{models} de
PostgreSQL y sube el archivo Python a MinIO; el
\texttt{model\_loader} de \texttt{simlab/model\_loader.py} descubre la
nueva entrada en la siguiente sesión y carga el archivo desde MinIO al
instanciar el modelo. Para añadir una nueva herramienta al Orchestrator
basta con registrar su esquema en el \textit{tool registry} de
\texttt{simlab/orchestrator.py} junto con su implementación; para
añadir un agente especializado, se sigue el mismo patrón que los
existentes, creando un nuevo módulo en \texttt{simlab/} y registrando
su \textit{tool call} en el Orchestrator. Los cambios sobre la capa de
persistencia compartida se incorporan mediante migraciones Alembic en
\texttt{shared/migrations/versions/}.
